\section{Elicitação de Requisitos}

\subsection{Reunião de grupo}
Foi utilizada uma discussão colaborativa entre os integrantes do grupo por mensagens no WhatsApp. A coleta ocorreu principalmente entre 29/09/2026, às 18:34, e 30/09/2026, às 18:58, em ambiente remoto. Participaram Pedro Luca Vaz Matias, Beatriz Montenegro e Frankley Kaiky M. Silva. A técnica foi escolhida porque o projeto ainda estava na fase de definição do tema e os integrantes precisavam alinhar a necessidade atendida, os atores e os riscos iniciais.

Como resultados, o grupo definiu:
\begin{itemize}
    \item O escopo principal do sistema, focando em um CRM ambulatorial (CliniCore) para gestão de relacionamento e agendamentos, descartando a complexidade de um prontuário eletrônico hospitalar completo.
    \item Os atores principais que interagem com o sistema: Administrador da Clínica, Recepcionista, Profissional de Saúde e Paciente.
    \item A necessidade de uma interface de interoperabilidade baseada no padrão HL7 FHIR.
\end{itemize}

\subsection{Análise Documental e Benchmarking}
Esta técnica consistiu na pesquisa de documentações técnicas e na avaliação de funcionalidades presentes em sistemas de gestão de clínicas já estabelecidos no mercado. A coleta ocorreu de forma assíncrona entre 30/09/2026 e 01/10/2026, em ambiente remoto, conduzida pelos mesmos integrantes do grupo. A técnica foi escolhida por ser altamente viável na ausência de stakeholders reais, permitindo que a equipe compreendesse os fluxos de trabalho padrão de uma clínica médica e os requisitos técnicos de interoperabilidade sem depender do contato direto com usuários finais.

\subsection{Considerações Finais}
A elicitação confirmou a necessidade de um sistema que separe a complexidade de um PEP robusto da agilidade exigida para a gestão do relacionamento e da agenda, enfatizando a necessidade de uma base de dados que seja facilmente interpretável por outros sistemas usando FHIR.